%!TeX root=../tese.tex
%("dica" para o editor de texto: este arquivo é parte de um documento maior)
% para saber mais: https://tex.stackexchange.com/q/78101

\chapter{Referencial Teórico}

\section{Desenvolvimento distribuído e organização do kernel Linux}

O desenvolvimento do kernel Linux constitui um caso representativo de produção distribuída de software livre em grande escala. Em contraste com projetos organizados em torno de uma equipe centralizada, o Linux adota uma estrutura hierárquica na qual autores de \textit{patches}, revisores e mantenedores de subsistemas exercem responsabilidades distintas. No nível inferior dessa hierarquia encontram-se os contribuidores, que propõem correções, novas funcionalidades e melhorias no código. Os mantenedores avaliam essas contribuições, decidem sobre sua aceitação e, quando apropriado, encaminham-nas para níveis superiores até que possam ser incorporadas ao repositório principal, mantido por Linus Torvalds (\cite{ferreira2019maintainers, hoang2021patchnet}).

Esse modelo combina duas infraestruturas complementares. As listas de \textit{e-mail} sustentam a submissão, a discussão e a revisão dos \textit{patches}, enquanto os repositórios Git registram as alterações aceitas e integradas ao código-fonte. Essa separação reflete a própria divisão entre revisão e integração. A primeira fase concentra-se na correção técnica, na utilidade e na conformidade da alteração com as práticas do projeto; a segunda envolve sua incorporação às árvores dos mantenedores e, posteriormente, à árvore principal. Consequentemente, uma proposta pode ser aprovada durante a revisão e, ainda assim, enfrentar conflitos ou outras dificuldades durante a integração (\cite{jiang2013patch}).

O processo pode iniciar-se com uma solicitação de comentários, denominada \textit{Request for Comments} (RFC), por meio da qual o desenvolvedor apresenta uma ideia de correção ou funcionalidade. Após considerar as manifestações recebidas, o autor implementa a alteração e envia o \textit{patch} ou uma série de \textit{patches} à lista do subsistema pertinente e, quando necessário, à lista geral do Linux. Especialistas e mantenedores analisam a proposta e podem aceitá-la, rejeitá-la ou solicitar modificações. Nesse último caso, o desenvolvedor produz uma nova versão e reinicia o ciclo de revisão. Quando a versão final é aprovada, o mantenedor a incorpora à sua árvore Git e a encaminha aos níveis superiores da hierarquia (\cite{ferreira2019maintainers, jiang2013patch}).

A integração à árvore de um mantenedor, portanto, não significa necessariamente que a alteração fará parte de uma versão oficial. Outros desenvolvedores podem realizar testes adicionais, identificar conflitos ou apontar problemas que impeçam a continuidade do processo. Se a contribuição superar essas etapas, ela poderá alcançar o repositório de Linus Torvalds, geralmente durante a \textit{merge window}, período de aproximadamente duas semanas subsequente ao lançamento de uma versão. Depois desse intervalo, a incorporação de novas funcionalidades é restringida, embora correções menores ainda possam ser admitidas até o lançamento seguinte (\cite{jiang2013patch}). 

Essa organização também se relaciona ao modelo de versões do Linux. A \textit{mainline} recebe novas funcionalidades e correções, ao passo que as versões estáveis recebem, em regra, correções previamente integradas à \textit{mainline}. Assim, um \textit{patch} destinado a uma versão estável deve primeiro ser incorporado ao repositório principal e, posteriormente, ser incorporado retroativamente para as versões afetadas. A identificação dessas alterações depende de critérios como impacto para o usuário, tamanho e correção aparente, mas sua aplicação envolve julgamento técnico e pode variar entre mantenedores (\cite{hoang2021patchnet}).

\section{O \textit{patch} como artefato técnico e comunicacional}

No processo do Linux, um \textit{patch} não se limita às linhas de código modificadas. Ele combina uma mensagem em linguagem natural e um diff que representa as exclusões e adições realizadas. O diff pode abranger vários arquivos, cada qual composto por um ou mais \textit{hunks}, que agrupam alterações localizadas. Essa estrutura torna o \textit{patch} um artefato híbrido: a mensagem explica a motivação e o propósito da contribuição, enquanto o código materializa tecnicamente a modificação (\cite{hoang2021patchnet}). 

A mensagem auxilia os revisores a compreender o problema, a justificativa da solução e seus possíveis efeitos. Por sua vez, as alterações no código fornecem evidências necessárias para verificar se a implementação corresponde à descrição. Essa complementaridade é relevante porque mensagens semelhantes podem acompanhar modificações distintas e, inversamente, alterações logicamente equivalentes podem ser explicadas com palavras diferentes. \cite{hoang2021patchnet} demonstram que tanto a mensagem quanto o código contribuem para a classificação de patches: o desempenho do PatchNet diminuiu quando qualquer uma dessas fontes foi retirada do modelo. 

As discussões desencadeadas pela submissão acrescentam uma dimensão social ao artefato. Revisores e mantenedores questionam escolhas de implementação, sugerem alternativas, solicitam testes e avaliam a compatibilidade da proposta com outros componentes. Por conseguinte, a trajetória de um \textit{patch} não pode ser compreendida exclusivamente pelo conteúdo da versão inicial ou do \textit{commit} final. As respostas e versões intermediárias registram parte do conhecimento produzido durante a revisão e permitem observar como a alteração foi negociada e transformada.

\cite{ferreira2019maintainers} caracterizam esse ambiente como sociotécnico, pois as decisões sobre o código são produzidas mediante interações entre pessoas com diferentes responsabilidades e conhecimentos. Em seu estudo, os autores analisaram 15.166 \textit{threads} que continham a participação de pelo menos um mantenedor. A necessidade de agrupar as mensagens em conversações decorre do fato de que a revisão se distribui por sucessivas respostas, e não apenas pela mensagem inicial. Dessa maneira, a \textit{thread} constitui uma unidade relevante para representar o contexto comunicacional de uma contribuição. 

Além da mensagem e do diff, metadados como autoria, destinatários, datas, assunto, \textit{Message-ID} e \textit{trailers} podem fornecer informações sobre o processo. Marcadores como \textit{Reviewed-by}, \textit{Tested-by} e \textit{Signed-off-by} registram formas distintas de participação e validação. Embora não sejam suficientes para comprovar a integração, esses elementos ajudam a caracterizar a maturidade da proposta, identificar os participantes e reconstruir sua passagem pelas etapas de revisão.

\section{Aceitação, revisão e integração de patches}

A aceitação de uma contribuição depende de fatores técnicos e processuais. Um \textit{patch} pode não avançar porque implementa uma solução inadequada, não segue as diretrizes do projeto, apresenta conflitos com alterações concorrentes ou não recebe atenção suficiente. Também é possível que a proposta seja reformulada e submetida novamente, dificultando a distinção entre rejeição definitiva, abandono e continuidade em outra \textit{thread}. Por isso, a ausência de um vínculo com um \textit{commit} não deve ser interpretada automaticamente como rejeição.

\cite{jiang2013patch} investigaram as características associadas à aceitação e ao tempo de integração. No período analisado, aproximadamente 33\% dos \textit{patches} submetidos foram incorporados a uma versão oficial, e a maioria dos aceitos levou de três a seis meses para chegar ao lançamento. Os autores decompuseram esse intervalo em tempo de revisão, tempo de integração e tempo de preparação para a versão. Essa distinção evidencia que aceitar tecnicamente um \textit{patch} e integrá-lo à \textit{mainline} são fenômenos relacionados, porém não equivalentes. 

A experiência do desenvolvedor, a maturidade do \textit{patch} e o histórico de alterações do subsistema estiveram entre os fatores associados à aceitação. Contribuidores com alterações previamente incorporadas apresentaram maior probabilidade de aceitação, enquanto propostas submetidas em versões sucessivas e precedidas por maior participação na discussão também tenderam a obter resultados favoráveis. Para o tempo de revisão, mostraram-se relevantes o momento da submissão, a quantidade de subsistemas afetados, o número de revisores indicados e a experiência do autor (\cite{jiang2013patch}). 

Esses achados reforçam que a identificação da integração não deve depender apenas da semelhança textual entre a submissão e o \textit{commit}. Os sinais sociais e processuais ajudam a explicar a evolução da contribuição e podem complementar as evidências extraídas do código. Entretanto, tais sinais precisam ser empregados com cautela: a autoria pode ser representada de maneiras diferentes; os destinatários podem mudar entre versões; e os trailers podem ser acrescentados, removidos ou reorganizados antes da integração.

A literatura sobre revisões também mostra que a comunicação pode influenciar a continuidade do processo. O estilo direto das discussões do Linux e a diversidade cultural de seus participantes tornam as interações parte relevante do ambiente de desenvolvimento. Embora \cite{ferreira2019maintainers} não tenham identificado mudanças amplas e consistentes no sentimento dos mantenedores durante o período estudado, o trabalho demonstra que as listas registram dimensões humanas inseparáveis das decisões técnicas.


\section{Perda de rastreabilidade entre \textit{e-mails} e \textit{commits}}

A existência de duas infraestruturas produz uma descontinuidade de rastreabilidade. Nas listas, o identificador básico é o \textit{Message-ID}; no Git, cada estado consolidado é identificado pelo \textit{hash} do \textit{commit}. Como não existe obrigatoriamente uma referência persistente entre esses identificadores, a associação pode desaparecer quando o \textit{patch} é integrado. Esse problema é agravado porque o Git preserva o resultado incorporado, mas não necessariamente toda a discussão ou todas as versões que antecederam a integração (\cite{ramsauer2019list}).

Durante a revisão, um \textit{patch} pode ser corrigido, reorganizado, dividido ou combinado com outras alterações. Operações de \textit{squash}, \textit{cherry-picking} e \textit{rebase}, bem como a resolução de conflitos, podem alterar a estrutura e o conteúdo da contribuição. Consequentemente, a relação entre os artefatos pode assumir cardinalidades de um para um, um para muitos ou muitos para um, dificultando a associação entre as versões submetidas e o commit integrado (\cite{ramsauer2019list}).

\cite{jiang2013patch} lidaram com esse problema dividindo os \textit{patches} em unidades correspondentes às alterações de cada arquivo. Em vez de comparar o \textit{patch} integralmente, os autores normalizaram as linhas modificadas e calcularam \textit{checksums} por unidade. Essa granularidade buscou acomodar situações de \textit{cherry-picking}, nas quais apenas partes da proposta são integradas. A abordagem obteve elevada precisão na amostra avaliada, mas também revelou limitações, sobretudo diante de mudanças de identificadores ou versões distribuídas em \textit{threads} diferentes. 

De forma complementar, \cite{jiang2014tracing} demonstram que técnicas baseadas na comparação das linhas adicionadas e removidas permitem recuperar vínculos entre \textit{e-mails} e \textit{commits} que não são registrados explicitamente. Esse método é mais tolerante a alterações nas linhas de contexto do que a comparação integral do diff. Ainda assim, a semelhança do código não elimina ambiguidades, pois modificações curtas ou recorrentes podem aparecer em diferentes arquivos e períodos.

Os índices Git constituem outra evidência relevante. Quando os índices dos \textit{blobs} anterior e posterior coincidem entre a submissão e o \textit{commit}, há uma indicação forte de que ambos representam a mesma transformação. Contudo, a divergência não implica necessariamente que sejam alterações distintas, já que qualquer modificação intermediária da base pode produzir novos objetos Git. Desse modo, os índices são apropriados como evidência positiva de elevada força, mas não como regra absoluta de exclusão.


\section{Vinculação de registros e combinação de evidências}
Sob a perspectiva metodológica, associar \textit{e-mails} a \textit{commits} configura um problema de vinculação de registros. O objetivo consiste em determinar se dois registros provenientes de fontes distintas representam a mesma alteração de software, apesar de diferenças de estrutura, conteúdo e granularidade. Em situações exatas, identificadores ou conteúdos idênticos permitem uma decisão praticamente determinística. Nos demais casos, a associação depende da combinação de atributos parcialmente discriminativos.

As evidências podem ser agrupadas em cinco dimensões. A dimensão técnica compreende arquivos modificados, linhas adicionadas e removidas, \textit{hunks}, índices de\textit{ blobs} e medidas de similaridade ou contenção. A dimensão textual abrange assunto, mensagem e descrição da alteração. A dimensão autoral inclui nomes, endereços e signatários. A dimensão temporal considera a ordem e a distância entre a submissão, a revisão e a integração. Finalmente, a dimensão processual reúne \textit{threads}, versões, trailers e vínculos explícitos, como referências Link.

A relevância de combinar mensagem e código também encontra suporte em \cite{hoang2021patchnet}. Embora o PatchNet tenha sido desenvolvido para identificar \textit{patches} destinados às versões estáveis, sua arquitetura demonstra que a representação conjunta dos componentes textuais e das mudanças no código supera representações baseadas em apenas uma dessas fontes. O modelo preserva a estrutura hierárquica do diff, distinguindo arquivos, \textit{hunks} e linhas adicionadas ou removidas, e combina essa representação com características aprendidas da mensagem. 

A aplicação dessa conclusão à vinculação entre \textit{e-mails} e \textit{commits} requer, contudo, uma distinção conceitual. A identificação de \textit{patches} estáveis é uma tarefa de classificação, enquanto a associação proposta neste trabalho é uma tarefa de correspondência entre pares. Ainda assim, ambos os problemas compartilham a necessidade de representar adequadamente o texto e o código, lidar com dados ruidosos e produzir medidas graduais em vez de decisões sustentadas por um único atributo.

A incerteza dos vínculos recomenda que o resultado seja expresso por níveis calibrados de confiança. Correspondências sustentadas por referência explícita, conteúdo técnico equivalente e metadados compatíveis podem receber confiança elevada. Quando apenas parte das evidências coincide, o vínculo pode ser classificado como provável. Por fim, evidências insuficientes ou contraditórias justificam uma categoria inconclusiva. Essa estrutura evita equiparar ausência de correspondência à rejeição e permite tratar separadamente erro de busca, transformação substancial e efetiva não integração.

\section{Escalabilidade, automação e avaliação}

A aplicação dessa vinculação ao histórico do Linux introduz um desafio de escala. A comparação de cada mensagem com todos os \textit{commits} produziria um número excessivo de pares. Por isso, a solução deve separar a geração de candidatos da classificação final. Na primeira etapa, critérios de menor custo, como janela temporal, autoria compatível e arquivos em comum, reduzem o universo de busca. Na segunda, medidas mais detalhadas avaliam o código, o texto e o contexto de revisão.

Essa arquitetura exige equilibrar eficiência e cobertura. Filtros excessivamente restritivos reduzem o custo computacional, mas podem eliminar correspondências verdadeiras antes da classificação. Em contrapartida, filtros muito amplos preservam a cobertura, porém transferem uma quantidade impraticável de candidatos para a etapa posterior. A escolha deve, portanto, ser avaliada não apenas pela precisão final, mas também pela revocação da geração de candidatos, pelo custo computacional e pela proporção de casos inconclusivos.

A automação do processo encontra respaldo em pesquisas que aplicam aprendizado de máquina à análise de \textit{patches}. O PatchNet, por exemplo, foi avaliado em 82.403 alterações e alcançou precisão média de 0,839, revocação de 0,907 e medida F1 de 0,871 para identificação de patches estáveis. Os resultados indicam que representações aprendidas conjuntamente a partir da mensagem e do código podem superar abordagens baseadas em palavras-chave ou atributos elaborados manualmente (\cite{hoang2021patchnet}). 

Contudo, a automação não elimina a necessidade de validação humana. O histórico utilizado como referência pode conter omissões, rótulos ruidosos e associações incompletas. Além disso, decisões sobre equivalência lógica podem exigir conhecimento do subsistema e interpretação da finalidade do código. Estudos recentes continuam apontando que, apesar do uso de ferramentas automáticas, a revisão de \textit{patches} do Linux permanece dependente de esforço humano especializado e enfrenta dificuldades para acompanhar o volume de submissões (\cite{lin2026reliable}). 

Nesse contexto, uma amostra de pares deve ser classificada independentemente por dois ou mais revisores, com mensuração da concordância antes da resolução das divergências. Essa amostra pode funcionar como verdade de referência para avaliar precisão, revocação, cobertura e calibração dos níveis de confiança. A revisão por consenso, realizada posteriormente, permite resolver desacordos sem ocultar a variabilidade inicial dos julgamentos.

Em síntese, a literatura estabelece três fundamentos para a pesquisa proposta. Primeiro, o processo hierárquico e distribuído do Linux separa comunicação, revisão e integração em diferentes infraestruturas. Segundo a evolução de um \textit{patch} pode modificar seu texto, seu código e sua granularidade, eliminando vínculos explícitos com o \textit{commit} final. Terceiro, a recuperação dessa trajetória requer uma abordagem escalável e multifatorial, capaz de combinar evidências técnicas, textuais, temporais e processuais. A lacuna situa-se, portanto, na construção de um método reproduzível que não apenas associe os artefatos, mas também mensure e calibre a confiança atribuída a cada vínculo.
